\clearpage
\section{Implementacja}

\subsection{Dodawanie i usuwanie węzłów}

Metoda \kod{addNode()} zapamiętuje rozmiar wektora i wywołuje
rekurencyjną metodę \kod{insert()}. Wzrost rozmiaru oznacza
dodanie nowego węzła i wynik \kod{true}. Wartość już obecna
w drzewie kończy operację wynikiem \kod{false}. Po powrocie
z wywołania rekurencyjnego rodzic jest ponownie odczytywany
przez indeks. Dzięki temu dodanie elementu pozostaje poprawne
także po przeniesieniu bufora wektora. Pokazuje to
listing~\ref{lst:wstawianie}.

\listingpliku[firstline=25,lastline=34,firstnumber=25,
  caption={Aktualizacja połączenia po wstawieniu do poddrzewa},
  label={lst:wstawianie}]{kod/BST.cpp}

Metoda \kod{removeNode()} wyszukuje iteracyjnie wartość i zapamiętuje
indeks rodzica. W przypadku dwojga dzieci szuka następnika, schodząc
najpierw w prawo, a następnie maksymalnie w lewo. Kopiuje jego
wartość do usuwanego węzła, a fizycznie usuwa następnika. Następnik
może mieć prawe dziecko, które przejmuje jego dotychczasowe miejsce.
W przypadku usuwania korzenia z najwyżej jednym dzieckiem
aktualizowane jest bezpośrednio pole \kod{rootIndex}.

Metoda usuwa pozycję z wektora dopiero po przełączeniu połączeń.
Potem zmniejsza wszystkie indeksy większe od usuniętej pozycji.
Zwraca \kod{true} po usunięciu węzła i \kod{false}, jeśli klucza
nie znaleziono. Metoda \kod{removeTree()} wywołuje \kod{tree.clear()}
i ustawia \kod{rootIndex} na $-1$. Wektor zachowuje pojemność do
ponownego wykorzystania.

\Needspace{7\baselineskip}
\subsection{Wyszukiwanie, przejścia i rysunek}

Metoda \kod{findElement()} zwraca informację o obecności klucza.
Używana przez menu metoda \kod{findPath()} dodatkowo zapisuje
wartości napotkanych węzłów. Znaleziona ścieżka obejmuje korzeń
i szukany węzeł. Gdy wyszukiwanie dojdzie do indeksu $-1$,
metoda zwraca pusty wektor. Jeżeli szukanym kluczem jest korzeń,
wynik zawiera jedną wartość.

\kod{preorderValues()} używa wektora indeksów jako stosu.
Pobiera węzeł, dopisuje jego wartość do wyniku, a następnie
odkłada na stos najpierw prawe, potem lewe dziecko. Dzięki
temu lewe poddrzewo jest odwiedzane wcześniej. Z wyniku tej
metody korzystają przejście preorder i oba tryby zapisu.
Inorder odwiedza lewe poddrzewo, węzeł i prawe poddrzewo,
uzyskując wartości rosnąco. Postorder wypisuje węzeł po obu
poddrzewach. Te dwa przejścia są rekurencyjne. Metoda
\kod{show()} jest osobną funkcją klasy i odczytuje wartości
w fizycznej kolejności wektora.

\kod{printTree()} wywołuje rekurencyjną metodę \kod{drawBranch()},
która tworzy etykiety z wartości liczbowych, przydziela im
miejsce w kolejności inorder i buduje wiersze rysunku. Między
poziomami drzewa umieszcza wiersze z łącznikami. Dla węzła
z dwojgiem dzieci pozycja środka etykiety jest wyznaczana
na podstawie pozycji obu dzieci. Dla pustego drzewa
\kod{printTree()} wypisuje komunikat \kod{Drzewo jest puste.}

\subsection{Sprawdzanie danych podczas odczytu}

\kod{loadBinary()} zaczyna od odczytania licznika, a następnie
wstawia kolejne wartości do nowego obiektu \kod{BST}. Wykrywa
brak wymaganych danych, duplikaty i nadmiarowe bajty.
Dopiero po sprawdzeniu końca pliku przenosi wynik do drzewa
używanego w menu. Poprawny zapis pustego drzewa zawiera
sam licznik o wartości zero. Wczytanie takiego pliku
opróżnia bieżące drzewo.

\kod{loadText()} zaczyna od kopii istniejącego drzewa. Czyta
tokeny oddzielone białymi znakami i sprawdza, czy każdy mieści
się w zakresie \kod{int} i składa się w całości z zapisu
liczby całkowitej. Duplikaty są pomijane przez \kod{addNode()}.
Pusty plik tekstowy zostaje odczytany poprawnie i nie dodaje
żadnych wartości. Przy błędzie odczytu albo błędnym tokenie
bieżące drzewo zachowuje swój stan także wtedy, gdy wcześniej
przetworzono już część pliku. Metody zapisu zwracają błąd
przy problemie z otwarciem, zapisem lub zamknięciem pliku.
Oba zapisy otwierają plik w trybie nadpisywania, więc błąd po
otwarciu może pozostawić plik częściowy.

\Needspace{12\baselineskip}
\subsection{Sprawdzenie działania}
\label{sec:testy}

Kod skompilowano w trybie C++14. Do pustego drzewa wprowadzono
kolejno liczby 8, 3, 10, 1, 6, 14, 4, 7 i 13. Listing~\ref{lst:przejscia}
zawiera wiersze z wynikami trzech przejść odczytane z konsoli.

\begin{lstlisting}[style=terminal,caption={Wyniki trzech przejść drzewa},label={lst:przejscia}]
Preorder: 8 3 1 6 4 7 10 14 13
Inorder: 1 3 4 6 7 8 10 13 14
Postorder: 1 4 7 6 3 13 14 10 8
\end{lstlisting}

Tabela~\ref{tab:proby} zestawia pozostałe próby wykonane z menu.
Próby importu wykonano z drzewa złożonego z węzła o wartości 8.

\begin{tabela}{Próby wykonane w menu}{Próby wykonane w menu}
              {\label{tab:proby}}
\small
\begin{tabularx}{\linewidth}{|L{4.8cm}|Y|}
\hline
\textbf{Próba} & \textbf{Zaobserwowany wynik}\\ \hline
Ponowne dodanie liczby 6 & Program zgłosił duplikat.\\ \hline
Wyszukanie 7 i 42 & Ścieżka do 7 to $8\to3\to6\to7$.
Program zgłosił brak 42.\\ \hline
Usunięcie korzenia 8 & Następnikiem został 10. Przejście inorder dało
$1,3,4,6,7,10,13,14$.\\ \hline
Zapis tekstowy i binarny & Plik tekstowy zawierał dziewięć wartości
w kolejności preorder. Odczyt zapisu binarnego po opróżnieniu drzewa
przywrócił ten sam wynik preorder i ten sam rysunek.\\ \hline
Odczyt obciętego pliku binarnego & Program zgłosił błąd,
a drzewo zachowało dotychczasową zawartość.\\ \hline
Import tekstu \kod{-4 8 9} &
Przejście inorder dało $-4,8,9$. Powtórzona wartość 8 wystąpiła raz.\\ \hline
Import tekstu \kod{-4 9x} &
Program zgłosił błąd. Drzewo nadal zawierało tylko wartość 8.\\ \hline
\end{tabularx}
\end{tabela}
